\subsubsection{紧性性质}\label{subsec:sobolev-compactness-properties}

本小节证明涉及 Sobolev 空间的某些嵌入是紧的. 这在应用中往往很有用, 因为它表明某些序列具有收敛子序列; 这种子序列的极限通常就是所求的解. 紧性也有助于定义 Sobolev 空间的某些等价范数, 本小节末尾就将这样做.

先定义紧嵌入.

\begin{Definition}
\label{def:compact-imbedding}
设 $X,Y$ 为 Banach 空间, 且 $X$ 连续嵌入 $Y$. 若 $X$ 中的单位球在 $Y$ 中预紧, 或等价地, $X$ 中每个有界序列都有在 $Y$ 中收敛的子序列, 则称 $X$ \textbf{紧嵌入} $Y$, 记为\[X\overset{c}{\hookrightarrow}Y\]
\end{Definition}

在确定 Sobolev 空间的紧嵌入之前, 先建立一个引理.

\begin{Lemma}
\label{lem:sobolev-holder-continuity}
设 $s>m/2$ 且 $0<\alpha<\min(s-m/2,1)$. 那么存在常数 $C$, 使对每个 $\bm{x},\bm{y}\in\mathbb{R}^m$ 及 $u\in H^s(\mathbb{R}^m)$ 有
\begin{equation}\label{eq:sobolev-holder-continuity-estimate}
|u(\bm{x})-u(\bm{y})|\leq C\|u\|_{s,2}|\bm{x}-\bm{y}|^{\alpha}
\end{equation}

\end{Lemma}

\begingroup
\postdisplaypenalty=10000
\renewcommand{\endProof}{\par\nobreak\hfill$\blacksquare$\par\vspace{0.1cm}}
\begin{Proof}
在整个证明中, 用 $C$ 表示与 $\bm{x}$, $\bm{y}$ 及 $u$ 无关的一般常数. 首先注意, 对任意 $\alpha\in(0,1)$,
\begin{equation}\label{eq:fourier-exponential-holder-difference}
|\exp(i\bm{\xi}\bm{x})-\exp(i\bm{\xi}\bm{y})|\leq C|\bm{x}-\bm{y}|^{\alpha}|\bm{\xi}|^{\alpha}.
\end{equation}
用 $\hat u$ 表示 $u$ 的 Fourier 变换, 则有
\begin{equation}\label{eq:sobolev-holder-fourier-estimate}
\begin{aligned}
|u(\bm{x})-u(\bm{y})|&\leq(2\pi)^{-m/2}\int_{\mathbb{R}^m}|\exp(i\bm{\xi}\bm{x})-\exp(i\bm{\xi}\bm{y})||\hat u(\bm{\xi})|\,\mathrm{d}\bm{\xi}\\
&\leq C|\bm{x}-\bm{y}|^{\alpha}\int_{\mathbb{R}^m}|\hat u(\bm{\xi})||\bm{\xi}|^{\alpha}\,\mathrm{d}\bm{\xi}\\
&\leq C|\bm{x}-\bm{y}|^{\alpha}\|(1+|\bm{\xi}|^s)\hat u\|_2\||\bm{\xi}|^{\alpha}/(1+|\bm{\xi}|^s)\|_2\\
&\leq C|\bm{x}-\bm{y}|^{\alpha}\|u\|_{s,2}.
\end{aligned}
\end{equation}

\end{Proof}
\endgroup

现在证明一个最基本的紧嵌入定理.

\begin{Theorem}
\label{thm:sobolev-continuous-target-compact-embedding}
设 $k>m/2$, $\Omega$ 有界且具有 $k$ 阶延拓性质. 那么\[H^k(\Omega)\overset{c}{\hookrightarrow}C_b(\overline\Omega),\]即 $H^k(\Omega)$ 紧嵌入 $C_b(\overline\Omega)$.
\end{Theorem}

\begin{Proof}
由 Sobolev 嵌入定理 (Theorem~\ref{thm:sobolev-full-space-continuous-embedding}), Lemma~\ref{lem:sobolev-holder-continuity} 及 Arzela-Ascoli 定理 (Theorem~\ref{thm:arzela-ascoli}) 立即得到本定理.
\end{Proof}

\Needspace{4\baselineskip}
\begin{Example}
\label{ex:unbounded-domain-translated-functions-noncompactness}
$\Omega$ 有界的假设是必不可少的. 为说明这一点, 令 $\Omega=\mathbb{R}^m$, 并考察平移函数序列\[u_n(\bm{x})=u(\bm{x}+n\bm{e}),\]其中 $\bm{e}$ 为任意固定单位向量, $u$ 为 $H^k(\mathbb{R}^m)$ 的任意非零元素. 当然, 该序列在 $H^k(\mathbb{R}^m)$ 中有界. 然而, 该序列不可能有一致收敛子序列. (为证明这一点, 需要由无穷远处的行为说明, 任何这样确实一致收敛的序列都不可能属于 $H^k(\mathbb{R}^m)$.)
\end{Example}
